% Options for packages loaded elsewhere
% Options for packages loaded elsewhere
\PassOptionsToPackage{unicode}{hyperref}
\PassOptionsToPackage{hyphens}{url}
\PassOptionsToPackage{dvipsnames,svgnames,x11names}{xcolor}
%
\documentclass[
  russian,
]{article}
\usepackage{xcolor}
\usepackage{amsmath,amssymb}
\setcounter{secnumdepth}{-\maxdimen} % remove section numbering
\usepackage{iftex}
\ifPDFTeX
  \usepackage[T1]{fontenc}
  \usepackage[utf8]{inputenc}
  \usepackage{textcomp} % provide euro and other symbols
\else % if luatex or xetex
  \usepackage{unicode-math} % this also loads fontspec
  \defaultfontfeatures{Scale=MatchLowercase}
  \defaultfontfeatures[\rmfamily]{Ligatures=TeX,Scale=1}
\fi
\usepackage{lmodern}
\ifPDFTeX\else
  % xetex/luatex font selection
  \setmainfont[]{DejaVu Serif}
  \setsansfont[]{DejaVu Sans}
  \setmonofont[]{DejaVu Sans Mono}
\fi
% Use upquote if available, for straight quotes in verbatim environments
\IfFileExists{upquote.sty}{\usepackage{upquote}}{}
\IfFileExists{microtype.sty}{% use microtype if available
  \usepackage[]{microtype}
  \UseMicrotypeSet[protrusion]{basicmath} % disable protrusion for tt fonts
}{}
\makeatletter
\@ifundefined{KOMAClassName}{% if non-KOMA class
  \IfFileExists{parskip.sty}{%
    \usepackage{parskip}
  }{% else
    \setlength{\parindent}{0pt}
    \setlength{\parskip}{6pt plus 2pt minus 1pt}}
}{% if KOMA class
  \KOMAoptions{parskip=half}}
\makeatother
% Make \paragraph and \subparagraph free-standing
\makeatletter
\ifx\paragraph\undefined\else
  \let\oldparagraph\paragraph
  \renewcommand{\paragraph}{
    \@ifstar
      \xxxParagraphStar
      \xxxParagraphNoStar
  }
  \newcommand{\xxxParagraphStar}[1]{\oldparagraph*{#1}\mbox{}}
  \newcommand{\xxxParagraphNoStar}[1]{\oldparagraph{#1}\mbox{}}
\fi
\ifx\subparagraph\undefined\else
  \let\oldsubparagraph\subparagraph
  \renewcommand{\subparagraph}{
    \@ifstar
      \xxxSubParagraphStar
      \xxxSubParagraphNoStar
  }
  \newcommand{\xxxSubParagraphStar}[1]{\oldsubparagraph*{#1}\mbox{}}
  \newcommand{\xxxSubParagraphNoStar}[1]{\oldsubparagraph{#1}\mbox{}}
\fi
\makeatother

\usepackage{color}
\usepackage{fancyvrb}
\newcommand{\VerbBar}{|}
\newcommand{\VERB}{\Verb[commandchars=\\\{\}]}
\DefineVerbatimEnvironment{Highlighting}{Verbatim}{commandchars=\\\{\}}
% Add ',fontsize=\small' for more characters per line
\usepackage{framed}
\definecolor{shadecolor}{RGB}{241,243,245}
\newenvironment{Shaded}{\begin{snugshade}}{\end{snugshade}}
\newcommand{\AlertTok}[1]{\textcolor[rgb]{0.68,0.00,0.00}{#1}}
\newcommand{\AnnotationTok}[1]{\textcolor[rgb]{0.37,0.37,0.37}{#1}}
\newcommand{\AttributeTok}[1]{\textcolor[rgb]{0.40,0.45,0.13}{#1}}
\newcommand{\BaseNTok}[1]{\textcolor[rgb]{0.68,0.00,0.00}{#1}}
\newcommand{\BuiltInTok}[1]{\textcolor[rgb]{0.00,0.23,0.31}{#1}}
\newcommand{\CharTok}[1]{\textcolor[rgb]{0.13,0.47,0.30}{#1}}
\newcommand{\CommentTok}[1]{\textcolor[rgb]{0.37,0.37,0.37}{#1}}
\newcommand{\CommentVarTok}[1]{\textcolor[rgb]{0.37,0.37,0.37}{\textit{#1}}}
\newcommand{\ConstantTok}[1]{\textcolor[rgb]{0.56,0.35,0.01}{#1}}
\newcommand{\ControlFlowTok}[1]{\textcolor[rgb]{0.00,0.23,0.31}{\textbf{#1}}}
\newcommand{\DataTypeTok}[1]{\textcolor[rgb]{0.68,0.00,0.00}{#1}}
\newcommand{\DecValTok}[1]{\textcolor[rgb]{0.68,0.00,0.00}{#1}}
\newcommand{\DocumentationTok}[1]{\textcolor[rgb]{0.37,0.37,0.37}{\textit{#1}}}
\newcommand{\ErrorTok}[1]{\textcolor[rgb]{0.68,0.00,0.00}{#1}}
\newcommand{\ExtensionTok}[1]{\textcolor[rgb]{0.00,0.23,0.31}{#1}}
\newcommand{\FloatTok}[1]{\textcolor[rgb]{0.68,0.00,0.00}{#1}}
\newcommand{\FunctionTok}[1]{\textcolor[rgb]{0.28,0.35,0.67}{#1}}
\newcommand{\ImportTok}[1]{\textcolor[rgb]{0.00,0.46,0.62}{#1}}
\newcommand{\InformationTok}[1]{\textcolor[rgb]{0.37,0.37,0.37}{#1}}
\newcommand{\KeywordTok}[1]{\textcolor[rgb]{0.00,0.23,0.31}{\textbf{#1}}}
\newcommand{\NormalTok}[1]{\textcolor[rgb]{0.00,0.23,0.31}{#1}}
\newcommand{\OperatorTok}[1]{\textcolor[rgb]{0.37,0.37,0.37}{#1}}
\newcommand{\OtherTok}[1]{\textcolor[rgb]{0.00,0.23,0.31}{#1}}
\newcommand{\PreprocessorTok}[1]{\textcolor[rgb]{0.68,0.00,0.00}{#1}}
\newcommand{\RegionMarkerTok}[1]{\textcolor[rgb]{0.00,0.23,0.31}{#1}}
\newcommand{\SpecialCharTok}[1]{\textcolor[rgb]{0.37,0.37,0.37}{#1}}
\newcommand{\SpecialStringTok}[1]{\textcolor[rgb]{0.13,0.47,0.30}{#1}}
\newcommand{\StringTok}[1]{\textcolor[rgb]{0.13,0.47,0.30}{#1}}
\newcommand{\VariableTok}[1]{\textcolor[rgb]{0.07,0.07,0.07}{#1}}
\newcommand{\VerbatimStringTok}[1]{\textcolor[rgb]{0.13,0.47,0.30}{#1}}
\newcommand{\WarningTok}[1]{\textcolor[rgb]{0.37,0.37,0.37}{\textit{#1}}}

\usepackage{longtable,booktabs,array}
\usepackage{calc} % for calculating minipage widths
% Correct order of tables after \paragraph or \subparagraph
\usepackage{etoolbox}
\makeatletter
\patchcmd\longtable{\par}{\if@noskipsec\mbox{}\fi\par}{}{}
\makeatother
% Allow footnotes in longtable head/foot
\IfFileExists{footnotehyper.sty}{\usepackage{footnotehyper}}{\usepackage{footnote}}
\makesavenoteenv{longtable}
\usepackage{graphicx}
\makeatletter
\newsavebox\pandoc@box
\newcommand*\pandocbounded[1]{% scales image to fit in text height/width
  \sbox\pandoc@box{#1}%
  \Gscale@div\@tempa{\textheight}{\dimexpr\ht\pandoc@box+\dp\pandoc@box\relax}%
  \Gscale@div\@tempb{\linewidth}{\wd\pandoc@box}%
  \ifdim\@tempb\p@<\@tempa\p@\let\@tempa\@tempb\fi% select the smaller of both
  \ifdim\@tempa\p@<\p@\scalebox{\@tempa}{\usebox\pandoc@box}%
  \else\usebox{\pandoc@box}%
  \fi%
}
% Set default figure placement to htbp
\def\fps@figure{htbp}
\makeatother



\ifLuaTeX
\usepackage[bidi=basic,provide=*]{babel}
\else
\usepackage[bidi=default,provide=*]{babel}
\fi
\ifPDFTeX
\else
\babelfont{rm}[]{DejaVu Serif}
\fi
% get rid of language-specific shorthands (see #6817):
\let\LanguageShortHands\languageshorthands
\def\languageshorthands#1{}


\setlength{\emergencystretch}{3em} % prevent overfull lines

\providecommand{\tightlist}{%
  \setlength{\itemsep}{0pt}\setlength{\parskip}{0pt}}



 


\makeatletter
\@ifpackageloaded{caption}{}{\usepackage{caption}}
\AtBeginDocument{%
\ifdefined\contentsname
  \renewcommand*\contentsname{Содержание}
\else
  \newcommand\contentsname{Содержание}
\fi
\ifdefined\listfigurename
  \renewcommand*\listfigurename{Список Иллюстраций}
\else
  \newcommand\listfigurename{Список Иллюстраций}
\fi
\ifdefined\listtablename
  \renewcommand*\listtablename{Список Таблиц}
\else
  \newcommand\listtablename{Список Таблиц}
\fi
\ifdefined\figurename
  \renewcommand*\figurename{Рисунок}
\else
  \newcommand\figurename{Рисунок}
\fi
\ifdefined\tablename
  \renewcommand*\tablename{Таблица}
\else
  \newcommand\tablename{Таблица}
\fi
}
\@ifpackageloaded{float}{}{\usepackage{float}}
\floatstyle{ruled}
\@ifundefined{c@chapter}{\newfloat{codelisting}{h}{lop}}{\newfloat{codelisting}{h}{lop}[chapter]}
\floatname{codelisting}{Список}
\newcommand*\listoflistings{\listof{codelisting}{Список Каталогов}}
\makeatother
\makeatletter
\makeatother
\makeatletter
\@ifpackageloaded{caption}{}{\usepackage{caption}}
\@ifpackageloaded{subcaption}{}{\usepackage{subcaption}}
\makeatother
\usepackage{bookmark}
\IfFileExists{xurl.sty}{\usepackage{xurl}}{} % add URL line breaks if available
\urlstyle{same}
\hypersetup{
  pdftitle={Этапы реализации проекта: Личный сайт на GitHub Pages},
  pdfauthor={Сунь Шэнцзе},
  pdflang={ru},
  colorlinks=true,
  linkcolor={blue},
  filecolor={Maroon},
  citecolor={Blue},
  urlcolor={Blue},
  pdfcreator={LaTeX via pandoc}}


\title{Этапы реализации проекта: Личный сайт на GitHub Pages}
\author{Сунь Шэнцзе}
\date{2026-05-26}
\begin{document}
\maketitle


\section{1. Цель
работы}\label{ux446ux435ux43bux44c-ux440ux430ux431ux43eux442ux44b}

Реализовать поддержку двух языков (русского и английского) на
персональном сайте, разместить контент на обоих языках, создать пост о
прошедшей неделе и тематический пост на двух языках.

\section{2. Этапы выполнения
проекта}\label{ux44dux442ux430ux43fux44b-ux432ux44bux43fux43eux43bux43dux435ux43dux438ux44f-ux43fux440ux43eux435ux43aux442ux430}

\subsection{2.1. Настройка поддержки двух
языков}\label{ux43dux430ux441ux442ux440ux43eux439ux43aux430-ux43fux43eux434ux434ux435ux440ux436ux43aux438-ux434ux432ux443ux445-ux44fux437ux44bux43aux43eux432}

В файле \texttt{config/\_default/languages.yaml} настроены русская и
английская версии сайта.

\begin{Shaded}
\begin{Highlighting}[]
\FunctionTok{ru}\KeywordTok{:}
\AttributeTok{  }\FunctionTok{languageCode}\KeywordTok{:}\AttributeTok{ ru{-}RU}
\AttributeTok{  }\FunctionTok{languageName}\KeywordTok{:}\AttributeTok{ Русский}
\AttributeTok{  }\FunctionTok{contentDir}\KeywordTok{:}\AttributeTok{ content/ru}
\AttributeTok{  }\FunctionTok{title}\KeywordTok{:}\AttributeTok{ }\StringTok{"Сунь Шэнцзе | Персональный сайт"}
\AttributeTok{  }\FunctionTok{weight}\KeywordTok{:}\AttributeTok{ }\DecValTok{1}

\FunctionTok{en}\KeywordTok{:}
\AttributeTok{  }\FunctionTok{languageCode}\KeywordTok{:}\AttributeTok{ en{-}US}
\AttributeTok{  }\FunctionTok{languageName}\KeywordTok{:}\AttributeTok{ English}
\AttributeTok{  }\FunctionTok{contentDir}\KeywordTok{:}\AttributeTok{ content/en}
\AttributeTok{  }\FunctionTok{title}\KeywordTok{:}\AttributeTok{ }\StringTok{"Sun Shengjie | Personal website"}
\AttributeTok{  }\FunctionTok{weight}\KeywordTok{:}\AttributeTok{ }\DecValTok{2}
\end{Highlighting}
\end{Shaded}

\subsection{2.2. Создание структуры
каталогов}\label{ux441ux43eux437ux434ux430ux43dux438ux435-ux441ux442ux440ux443ux43aux442ux443ux440ux44b-ux43aux430ux442ux430ux43bux43eux433ux43eux432}

Созданы каталоги:

\begin{Shaded}
\begin{Highlighting}[]
\NormalTok{content/ru/}
\NormalTok{content/en/}
\end{Highlighting}
\end{Shaded}

Внутри были созданы подкаталоги:

\begin{itemize}
\tightlist
\item
  authors
\item
  blog
\item
  project
\item
  home
\end{itemize}

После этого существующий контент был переведён на английский язык.

\subsection{2.3. Адаптация авторской
информации}\label{ux430ux434ux430ux43fux442ux430ux446ux438ux44f-ux430ux432ux442ux43eux440ux441ux43aux43eux439-ux438ux43dux444ux43eux440ux43cux430ux446ux438ux438}

Для английской версии создан файл:

\begin{Shaded}
\begin{Highlighting}[]
\NormalTok{content/en/authors/me/\_index.md}
\end{Highlighting}
\end{Shaded}

Содержание файла:

\begin{Shaded}
\begin{Highlighting}[]
\CommentTok{{-}{-}{-}}
\AnnotationTok{title:}\CommentTok{ "Sun Shengjie"}
\AnnotationTok{role:}\CommentTok{ "student"}
\AnnotationTok{bio:}\CommentTok{ "RUDN student, interested in operating systems and computer architecture."}

\AnnotationTok{interests:}
\CommentTok{  {-} Operating Systems}
\CommentTok{  {-} Computer Architecture}

\AnnotationTok{education:}
\CommentTok{  courses:}
\CommentTok{    {-} course: "Bachelor\textquotesingle{}s degree (02.03.00)"}
\CommentTok{      institution: "RUDN University"}
\CommentTok{      year: 2025}
\CommentTok{{-}{-}{-}}
\end{Highlighting}
\end{Shaded}

\subsection{2.4. Создание постов на двух
языках}\label{ux441ux43eux437ux434ux430ux43dux438ux435-ux43fux43eux441ux442ux43eux432-ux43dux430-ux434ux432ux443ux445-ux44fux437ux44bux43aux430ux445}

\subsubsection{2.4.1. Пост о прошедшей
неделе}\label{ux43fux43eux441ux442-ux43e-ux43fux440ux43eux448ux435ux434ux448ux435ux439-ux43dux435ux434ux435ux43bux435}

Созданы файлы:

\begin{Shaded}
\begin{Highlighting}[]
\NormalTok{content/ru/blog/week4{-}bilingual/index.md}
\NormalTok{content/en/blog/week4{-}bilingual/index.md}
\end{Highlighting}
\end{Shaded}

Русская версия:

\begin{Shaded}
\begin{Highlighting}[]
\CommentTok{{-}{-}{-}}
\AnnotationTok{title:}\CommentTok{ "Четвёртая неделя: настройка двуязычного сайта"}
\AnnotationTok{date:}\CommentTok{ 2026{-}05{-}26}
\AnnotationTok{authors:}
\CommentTok{  {-} me}
\CommentTok{{-}{-}{-}}

\NormalTok{На этой неделе я настроил поддержку двух языков.}
\NormalTok{Сайт теперь доступен на русском и английском языках.}
\end{Highlighting}
\end{Shaded}

\pandocbounded{\includegraphics[keepaspectratio]{images/project_26_bilingual_post_ru.png}}

\emph{Рисунок 1: Русская версия поста}

\pandocbounded{\includegraphics[keepaspectratio]{images/project_27_bilingual_post_en.png}}

\emph{Рисунок 2: Английская версия поста}

\subsubsection{2.4.2. Тематический пост о
Markdown}\label{ux442ux435ux43cux430ux442ux438ux447ux435ux441ux43aux438ux439-ux43fux43eux441ux442-ux43e-markdown}

Созданы файлы:

\begin{Shaded}
\begin{Highlighting}[]
\NormalTok{content/ru/blog/markdown{-}bilingual/index.md}
\NormalTok{content/en/blog/markdown{-}bilingual/index.md}
\end{Highlighting}
\end{Shaded}

В постах описаны возможности языка разметки Markdown и особенности его
использования при создании документации и блогов.

\section{3.
Заключение}\label{ux437ux430ux43aux43bux44eux447ux435ux43dux438ux435}

В результате выполнения этапа:

\begin{itemize}
\tightlist
\item
  Настроена поддержка двух языков
\item
  Создана структура каталогов для мультиязычного контента
\item
  Переведена авторская информация
\item
  Созданы посты на русском и английском языках
\end{itemize}

Сайт успешно опубликован на GitHub Pages.

Адрес сайта:

\begin{Shaded}
\begin{Highlighting}[]
\NormalTok{https://oiopuppy.github.io}
\end{Highlighting}
\end{Shaded}





\end{document}
